\scanpage{77}
% lines 1–4, universal pair
Ceci posé,
\[
\bigl(\deleted{E}P_1*P_2,\quad P_1\mathbin{\unsure{*|\times}}P_2\xrightarrow{*}P_1*P_2\bigr)
\]
\deleted{\unsure{avec}} est une solution du pb universel défini par cette notion (pour $M,\Pi$ variables...).
% lines 5–14, compatible trivializations
Si on veut de plus décrire $P_1*P_2$ comme un torseur $\unsure{\chi|\gamma}$-trivialisé, on \deleted{définira} exigera, dans la notion de ``morphisme biaffine compatible avec les $\chi$-trivialisations'' (non seulement de $P_1,P_2$, mais aussi de $\Pi$), que l'on ait
\[
\pi_{\Pi}(\deleted{\unsure{\alpha}}f(x_1,x_2))=\pi_1(x_1)\otimes\pi_2(x_2)
\]
et on trouve maintenant une solution \unsure{de|du} ce problème universel pour le torseur $\chi$-trivialisé $P_1*P_2$.
% lines 15–20, morphisms
Mais on avait oublié d'expliciter la notion de morphisme de torseurs $\chi$-trivialisés,
\[
\begin{aligned}
P&\xrightarrow{f}P'\\
M&\xrightarrow{f_0}M'
\end{aligned}
\]
On exige
\[
\left\{\begin{aligned}
f(\overset{\substack{\Gamma P\\\rotatebox{-90}{$\in$}}}{x}+\overset{\substack{\Gamma M\\\rotatebox{-90}{$\in$}}}{u})&=f(x)+f_0(u)\\
\pi_{P'}(f(x))&=f_0\bigl(\pi_P(x)\bigr)
\end{aligned}\right.
\]
% archive sheet number: 76
